\documentclass[11pt]{article}

\usepackage[T1]{fontenc}
\usepackage[utf8]{inputenc}
\usepackage{lmodern}
\usepackage{amsmath}
\usepackage{amssymb}
\usepackage{microtype}
\usepackage{hyperref}
\usepackage{booktabs}

\title{Regulatory Compliance Metric}
\author{Alquimia AI}
\date{}

\begin{document}

\maketitle

\section{Regulatory Compliance Metric}

The Regulatory Compliance Metric evaluates whether an AI assistant's responses adhere to a corpus of regulatory documents. Rather than relying on an LLM judge, this metric grounds its evaluation in a retrieval-augmented pipeline: relevant passages are retrieved from the regulatory corpus and then scored against the assistant's response using a neural reranker. This design keeps evaluation transparent and auditable, producing verdicts that are directly traceable to specific regulatory passages.

\subsection{Motivation}

Deployed AI assistants in regulated domains---such as healthcare, finance, or legal services---must produce responses that align with governing regulations, policies, and compliance requirements. A mismatch between an assistant's output and a regulatory document may constitute a compliance violation, regardless of how fluent or helpful the response appears. The Regulatory Compliance Metric addresses this gap by treating the regulatory corpus as a ground truth and quantifying the degree to which each assistant response is supported or contradicted by that corpus.

\subsection{Retrieval-Augmented Evaluation Pipeline}

The evaluation pipeline proceeds in three stages: corpus ingestion, chunk retrieval, and contradiction checking.

\subsubsection{Corpus Ingestion and Chunking}

Regulatory documents are loaded via a \texttt{CorpusConnector} abstraction and chunked into fixed-size overlapping windows by the \texttt{DocumentRetriever}. Given a document of length $L$ characters, chunk size $c$, and overlap $o$, the $k$-th chunk spans characters $[k(c - o),\; k(c - o) + c]$. Consecutive whitespace normalization is applied before chunking to reduce noise. The total number of chunks across all documents constitutes the retrieval index.

All chunk texts are encoded into dense vectors by an \texttt{Embedder} instance. The \texttt{Embedder} interface exposes two encoding methods: \texttt{encode} for documents and \texttt{encode\_query} for queries, allowing implementations to apply instruction prefixes or asymmetric encoding strategies typical of bi-encoder retrieval models~\cite{reimers2019sentencebert}.

\subsubsection{Merged Chunk Retrieval}

For each interaction, the retriever issues two separate embedding lookups---one against the user query $Q$ and one against the assistant response $A$---and merges the results:

\begin{equation}
    \mathcal{C}(Q, A) = \text{deduplicate}\!\left(\text{retrieve}(Q) \cup \text{retrieve}(A)\right)
\end{equation}

where deduplication retains, for each unique (source, chunk\_index) pair, the chunk with the highest cosine similarity. This dual-query strategy ensures that chunks relevant to the question's subject matter and chunks topically aligned with the assistant's claims are both surfaced.

Cosine similarity is used as the retrieval scoring function. Let $\mathbf{q}$ be the query embedding and $\mathbf{d}_i$ the embedding of the $i$-th chunk; only chunks satisfying $\text{sim}(\mathbf{q}, \mathbf{d}_i) \geq \tau_s$ are retained, where $\tau_s$ is the configurable similarity threshold (default $\tau_s = 0.3$). The top-$k$ chunks by similarity are forwarded to the contradiction checker (default $k = 10$).

\subsubsection{Contradiction Checking via Reranking}

The retrieved chunks are re-scored by a \texttt{Reranker} that computes a relevance score $r_i \in [0, 1]$ for each chunk $\mathbf{c}_i$ with respect to the assistant response $A$. A cross-encoder architecture is typical here, as it jointly encodes the document-response pair rather than relying on independent embeddings~\cite{nogueira2019passage}.

Each chunk is assigned a binary verdict based on a configurable contradiction threshold $\tau_c$ (default $\tau_c = 0.6$):

\begin{equation}
    v_i = \begin{cases}
        \texttt{SUPPORTS}     & \text{if } r_i \geq \tau_c \\
        \texttt{CONTRADICTS}  & \text{if } r_i < \tau_c
    \end{cases}
\end{equation}

\subsection{Interaction-Level Verdict and Compliance Score}

Let $n_s$ and $n_c$ denote the number of supporting and contradicting chunks for a given interaction, respectively, and $n_t = n_s + n_c$ the total number of retrieved chunks.

If no chunks are retrieved ($n_t = 0$), the interaction is assigned an \texttt{IRRELEVANT} verdict, indicating that the regulatory corpus does not cover the topic of the query-response pair. Otherwise, the raw compliance score is:

\begin{equation}
    s = \frac{n_s}{n_t}
\end{equation}

The interaction-level verdict is then determined by comparing $s$ against a compliance threshold $\tau_p$ (default $\tau_p = 0.5$):

\begin{equation}
    \text{verdict} = \begin{cases}
        \texttt{NON\_COMPLIANT} & \text{if } n_c > 0 \text{ and } n_s = 0 \\
        \texttt{COMPLIANT}      & \text{if } s \geq \tau_p \\
        \texttt{NON\_COMPLIANT} & \text{otherwise}
    \end{cases}
\end{equation}

The first branch ensures that any interaction where every retrieved chunk contradicts the response is immediately classified as non-compliant, regardless of the threshold.

\subsection{Session-Level Aggregation}

Interaction-level compliance scores are aggregated to a session-level metric using a pluggable \texttt{StatisticalMode} interface that supports both frequentist and Bayesian paradigms. For a session with $n$ interactions producing scores $s_1, \ldots, s_n$ and weights $w_1, \ldots, w_n$:

\paragraph{Frequentist Mode.} Returns the weighted mean compliance score:

\begin{equation}
    \bar{s} = \sum_{i=1}^{n} w_i \cdot s_i
\end{equation}

\paragraph{Bayesian Mode.} Returns a bootstrapped posterior with credible intervals. Given the score vector $\mathbf{s} = (s_1, \ldots, s_n)$ and weight vector $\mathbf{w} = (w_1, \ldots, w_n)$, we draw $S$ Monte Carlo bootstrap samples (default $S = 5000$):

\begin{equation}
    \bar{s}^{(k)} = \frac{1}{n} \sum_{j=1}^{n} s_{I_j^{(k)}}, \quad I_j^{(k)} \sim \text{Categorical}(\mathbf{w}), \quad k = 1, \ldots, S
\end{equation}

The session-level compliance is then reported with a credible interval:

\begin{equation}
    \hat{s} = \mathbb{E}[\bar{s}^{(k)}], \quad \text{CI}_{1-\alpha} = \left[ Q_{\alpha/2}(\bar{s}^{(k)}), \; Q_{1-\alpha/2}(\bar{s}^{(k)}) \right]
\end{equation}

where $Q_p$ denotes the $p$-th quantile of the bootstrap distribution and $\alpha = 0.05$ for 95\% credible intervals.

The session-level verdict mirrors the interaction-level logic: if all interactions are \texttt{IRRELEVANT}, the session is \texttt{IRRELEVANT}; otherwise, the session is \texttt{COMPLIANT} if $\bar{s} \geq \tau_p$ and \texttt{NON\_COMPLIANT} otherwise.

\subsection{Data Representation}

Evaluation results are represented through a typed schema hierarchy:

\begin{enumerate}
    \item \textbf{RegulatoryChunk}: Stores a single retrieved chunk with its full provenance:
    \begin{itemize}
        \item \texttt{text}, \texttt{source}, \texttt{chunk\_index}: The chunk content and its location within the corpus
        \item \texttt{similarity}: Cosine similarity from embedding retrieval ($\in [0, 1]$)
        \item \texttt{reranker\_score}: Cross-encoder score against the assistant response ($\in [0, 1]$)
        \item \texttt{verdict}: \texttt{SUPPORTS} or \texttt{CONTRADICTS}
    \end{itemize}
    \item \textbf{RegulatoryInteraction}: Per-interaction result containing the compliance score, verdict, supporting and contradicting chunk counts, the full list of \texttt{RegulatoryChunk} records, and a human-readable \texttt{insight} string.
    \item \textbf{RegulatoryMetric}: Session-level aggregate containing:
    \begin{itemize}
        \item \texttt{session\_id} and \texttt{assistant\_id}: Session and assistant identifiers
        \item \texttt{n\_interactions}: Total number of interactions evaluated
        \item \texttt{compliance\_score}: Aggregated session compliance score $\bar{s}$
        \item \texttt{compliance\_score\_ci\_low}, \texttt{compliance\_score\_ci\_high}: Credible interval bounds (Bayesian mode only; \texttt{None} otherwise)
        \item \texttt{verdict}: Session-level \texttt{COMPLIANT}, \texttt{NON\_COMPLIANT}, or \texttt{IRRELEVANT}
        \item \texttt{total\_supporting\_chunks}, \texttt{total\_contradicting\_chunks}: Corpus-level evidence summary
        \item \texttt{interactions}: Full list of \texttt{RegulatoryInteraction} records for drill-down analysis
    \end{itemize}
\end{enumerate}

All schemas are implemented as Pydantic~\cite{pydantic2024} models with strict field validation (e.g., \texttt{compliance\_score: float = Field(ge=0, le=1)}), ensuring that scores remain within well-defined bounds throughout the evaluation pipeline.

\subsection{Architectural Considerations}

The regulatory pipeline follows the same dependency inversion and strategy patterns used throughout the framework. The \texttt{Embedder} and \texttt{Reranker} abstractions are defined as abstract base classes; concrete implementations are injected at construction time, decoupling the pipeline from any specific model provider. This allows practitioners to substitute, for example, a local sentence-transformer embedder with a cloud-hosted embedding API without modifying the evaluation logic.

The corpus is loaded lazily on the first call to \texttt{batch()}, ensuring that initialization cost is deferred until evaluation actually begins. Subsequent calls reuse the pre-built embedding index, so the $O(|D|)$ embedding cost is paid only once per \texttt{Regulatory} instance regardless of the number of sessions evaluated.

\section{Proposed Experiments}

A key open question in deploying the Regulatory Compliance Metric is how the choice of embedding model and reranker jointly affect the quality of retrieval and contradiction detection. We propose a benchmark that systematically evaluates these components against a fixed regulatory target.

\subsection{Experimental Setup}

We select the EU AI Act~\cite{euaiact2024} as the target regulatory corpus, given its breadth, public availability, and practical relevance to AI system evaluation. A human-annotated test set is constructed as follows:

\begin{enumerate}
    \item A set of $N$ query-response pairs is sampled from a conversational AI system operating in a regulated domain.
    \item For each pair, domain experts label the subset of EU AI Act articles that are normatively relevant to the assistant's response.
    \item Each labeled article is further annotated with a binary verdict---\texttt{SUPPORTS} or \texttt{CONTRADICTS}---indicating whether the article's normative content is consistent with the response.
\end{enumerate}

This annotated set serves as the ground truth for evaluating both retrieval quality and contradiction detection accuracy.

\subsection{Embedding Model Comparison}

We evaluate a set of bi-encoder embedding models spanning different architectures, training objectives, and parameter counts. Candidate models include general-purpose sentence encoders (e.g., \texttt{all-mpnet-base-v2}), instruction-tuned variants (e.g., \texttt{e5-large-instruct}), and domain-adapted legal embedders where available. For each model, retrieval quality is measured using standard information retrieval metrics computed against the ground-truth relevant articles:

\begin{itemize}
    \item \textbf{Recall@$k$}: Fraction of ground-truth relevant chunks retrieved within the top-$k$ results. This is the primary metric, as missing a relevant article is a more critical failure than including an irrelevant one.
    \item \textbf{Mean Reciprocal Rank (MRR)}: Average reciprocal rank of the first relevant chunk across queries, measuring how highly relevant content is ranked.
    \item \textbf{Normalized Discounted Cumulative Gain (nDCG@$k$)}: Ranking quality accounting for graded relevance across positions.
\end{itemize}

\subsection{Reranker Comparison}

Given a fixed retrieval set (top-$k$ chunks from the best-performing embedding model), we evaluate a set of cross-encoder rerankers on their ability to correctly classify each chunk as \texttt{SUPPORTS} or \texttt{CONTRADICTS}. Candidate rerankers include general cross-encoders (e.g., \texttt{cross-encoder/ms-marco-MiniLM-L-6-v2}), larger rerankers (e.g., \texttt{bge-reranker-large}), and any available models fine-tuned on legal or regulatory text. Evaluation metrics are:

\begin{itemize}
    \item \textbf{Classification F1}: Macro-averaged F1 over the \texttt{SUPPORTS} / \texttt{CONTRADICTS} binary classification task, using the ground-truth annotations as labels.
    \item \textbf{Threshold sensitivity}: F1 as a function of the contradiction threshold $\tau_c \in [0, 1]$, identifying the operating point that maximises performance on the target corpus.
    \item \textbf{Precision-Recall curve}: Full curve across thresholds, allowing practitioners to trade off false positives (over-flagging compliant responses) against false negatives (missing genuine violations).
\end{itemize}

\subsection{Joint Evaluation}

Finally, we evaluate embedding-reranker combinations end-to-end using the session-level compliance score as the outcome variable. For each pair, the predicted session verdict (\texttt{COMPLIANT} / \texttt{NON\_COMPLIANT}) is compared against a human-assigned session verdict, and overall accuracy and Cohen's $\kappa$ are reported. This joint evaluation captures interaction effects: a strong reranker may compensate for weaker retrieval, and vice versa, so individual component scores do not necessarily predict end-to-end performance.

Results are expected to inform practical recommendations for corpus-specific model selection, and to surface whether domain-adapted models provide meaningful gains over general-purpose alternatives on regulatory text.

\bibliographystyle{plain}
\bibliography{references}

\end{document}